\SourceLeaf{429}%
\begin{center}10.\end{center}
\noindent\hspace*{1.2em}Ehe wir den eben gefundenen Ausdruck weiter umformen, wollen wir über denselben einige Bemerkungen machen.

Eine Normale auf die Curve in ihrer Ebene entspricht zweien Richtungen auf der Kugelfläche, je nachdem man jene auf die eine oder auf die andere Seite der Curve zieht; die eine Richtung, nach welcher die Curve \textls[100]{concav} ist, wird durch $\lambda'$ bezeichnet, die andere durch den gegenüberliegenden Punkt der Kugelfläche. Diese beiden Punkte, eben so wie $L$ und $\mathfrak{L}$, sind von $\lambda$ um $90^{0}$ entfernt und liegen also in einem grössten Kreise; und da auch $\mathfrak{L}$ von $\lambda$ um $90^{0}$ entfernt ist, so wird $\mathfrak{L}L=90^{0}-L\lambda'$ oder $=L\lambda'-90^{0}$ sein, und daher
\[
\cos L\lambda'=\pm\sin\mathfrak{L}L,
\]
wo $\sin\mathfrak{L}L$ nothwendig positiv ist. Da in unserer Analyse $r$ als positiv betrachtet ist, so wird das Zeichen von $\cos L\lambda'$ dem von
\[
Z(\xi\xi T+2\xi\eta U+\eta\eta V)
\]
gleich sein, und also ein positiver Werth von letzterer Grösse anzeigen, dass $L\lambda'$ unter $90^{0}$ ist, oder dass die Curve nach der Seite concav ist, wo die Projection auf die Ebene der Normale auf die krumme Fläche liegt, ein negativer Werth hingegen, dass die Curve nach dieser Seite convex ist. Wir können daher auch allgemein
\[
\frac1r=\frac{Z(\xi\xi T+2\xi\eta U+\eta\eta V)}{\sin\mathfrak{L}L}
\]
setzen, wenn wir den Krümmungshalbmesser im ersten Fall als positiv, im andern als negativ betrachten wollen. $\mathfrak{L}L$ ist nunmehro der Winkel, welchen unsere schneidende Ebene mit der die krumme Fläche berührenden Ebene macht, und man sieht, dass in verschiedenen schneidenden Ebenen, die durch denselben Punkt und dieselbe Tangente gelegt sind, die Krümmungshalbmesser dem Sinus der Neigung proportional sind. Wegen dieser einfachen Beziehung wollen wir uns hinfort auf den Fall einschränken, wo dieser Winkel ein rechter ist und also die schneidende Ebene durch die Normale auf die krumme
